% ============================================================
% analy 报告 — test13 多线程 CUDA C++ MultiGrid 测试工作
% 编译: xelatex 两遍；交付前 Overfull 计数必须为 0
% ============================================================
\documentclass[UTF8]{ctexart}
\usepackage[margin=2.5cm]{geometry}
\usepackage{amsmath}
\usepackage{microtype}
\usepackage{listings}
\usepackage{xcolor}
\usepackage{booktabs}
\usepackage{longtable}
\usepackage{xltabular}
\usepackage{tabularx}
\usepackage{hyperref}
\usepackage{fancyhdr}
\usepackage[most]{tcolorbox}
\usepackage{titlesec}

% ===== 色板 =====
\definecolor{accent}{HTML}{1F4E79}
\definecolor{evidence}{HTML}{1E8449}
\definecolor{reference}{HTML}{8C6D1F}
\definecolor{conclusion}{HTML}{8B1A1A}
\definecolor{codebg}{HTML}{F4F6F7}
\definecolor{struct}{HTML}{3B3B98}
\definecolor{relation}{HTML}{0E7C7B}
\definecolor{idea}{HTML}{B03A2E}
\definecolor{warn}{HTML}{D35400}
\definecolor{hl}{HTML}{FDF6E3}

\setlength{\emergencystretch}{3em}
\hypersetup{colorlinks=true,linkcolor=accent,urlcolor=relation}

\titleformat{\section}{\Large\bfseries\color{accent}}{\thesection}{1em}{}
\titleformat{\subsection}{\large\bfseries\color{struct}}{\thesubsection}{1em}{}
\titleformat{\subsubsection}{\normalsize\bfseries\color{struct!70!black}}{\thesubsubsection}{1em}{}

\newtcolorbox{conclbox}{breakable,colback=conclusion!6,colframe=conclusion,title=结论}
\newtcolorbox{refbox}{breakable,colback=reference!6,colframe=reference,title=参考背景}
\newtcolorbox{ideabox}{breakable,colback=idea!6,colframe=idea,title=项目思路}
\newtcolorbox{warnbox}{breakable,colback=warn!6,colframe=warn,title=局限与风险}
\newtcolorbox{keybox}{breakable,colback=hl,colframe=accent,title=要点}

\lstset{
  basicstyle=\ttfamily\footnotesize,
  backgroundcolor=\color{codebg},
  frame=single,
  language=python,
  firstnumber=auto,
  keywordstyle=\color{accent}\bfseries,
  commentstyle=\color{evidence!70!black},
  stringstyle=\color{struct},
  breaklines=true,
  breakatwhitespace=false,
  postbreak=\mbox{\textcolor{accent}{$\hookrightarrow$}\space},
  keepspaces=true,
  columns=flexible,
  numbers=left,numberstyle=\tiny\color{struct!60!black},
}

\title{test13 多线程 CUDA C++ MultiGrid 测试工作分析报告}
\author{opencode analy}
\date{2026-08-15}

\begin{document}
\maketitle

\begin{abstract}
本报告分析 \texttt{logs/test13} 测试套件及其支撑代码在 2026-08-15 的收尾工作：多线程
（一线程一卡）\texttt{MultiGpuMultigrid} 求解器的正确性验证、性能基准与参数扫描全流程
执行，以及测试脚本三处读取 bug 的修复。核心结论：正确性全部 PASS；加速比 gate=1.5
通过 11/16 配置，最佳 2.149（8x8x8x16, L3 r10）；小格子 2L 多线程加速 2.10 倍，
中/大格子加速比 $<1$ 为已知历史特性。三视角概览：主体结构（套件=单文件 main.py +
运行脚本 + 版本化产物目录）、各部分关系（main.py 子命令链 + h5py 数据流 + nullvec
缓存复用）、项目思路（真多线程并行、加速比语义、数据持久化纯 h5py）。
\end{abstract}

\tableofcontents

% ================= 1 仓库概览 =================
\section{仓库概览}
\subsection{分析范围}
本次分析聚焦 \texttt{logs/test13} 套件（816 行 main.py + run-local.sh + AGENTS.md +
14 项运行产物）及其核心支撑 \texttt{pyqcu/cuda/\_multi\_gpu.py}（534 行）；主题为
本次会话完成的工作：verify/bench/sweep/check/collect/mktable/plots/budget 全流程、
\texttt{\_multi\_gpu.py} 粗算子构建统一 C++ matvec 路径（前置未提交改动，经 verify
确认正确）、main.py 三处读取 bug 修复。

\subsection{规模统计与 git 信息}
\begin{table}[htbp]\centering
\begin{tabular}{ll}
\toprule
指标 & 实测值 \\
\midrule
套件代码文件 & main.py (816 行) + run-local.sh + AGENTS.md \\
运行产物 & 14 项（h5/tex/png）归档于 v202608151626/ \\
nullvec 缓存 & 8 个粗算子 h5（0.1--2.8 GB） \\
提交数 & 最新提交 eef6bbf (2026-08-15 04:31) \\
最近标签 & test12\_2 (2026-08-14) \\
\bottomrule
\end{tabular}
\caption{分析对象实测统计（数据来源: wc/git/ls 实测）}
\end{table}

% ================= 2 主体结构 =================
\section{主体结构}
\begin{table}[htbp]\centering
\begin{tabularx}{\textwidth}{llX}
\toprule
\textcolor{struct}{层次} & 路径 & \textcolor{struct}{职责} \\
\midrule
入口层 & logs/test13/main.py & 全部子命令（verify/clean/bench/sweep/check/\\
&& budget/collect/mktable/plots），--outdir 公共参数 \\
编排层 & logs/test13/run-local.sh & 全流程串联 + timeout 防卡壳 + tee 归档 \\
被测核心 & pyqcu/cuda/\_multi\_gpu.py & MultiGpuMultigrid：N 线程×卡并行求解 \\
支撑核心 & pyqcu/cuda/\_schur\_op.py & CudaSchurOp：每线程独立 C++ matvec 算子 \\
数据层 & logs/test13/*.h5 + nullvec\_cache/ & h5py 结果持久化与粗算子缓存 \\
文档层 & logs/test13/AGENTS.md, test13\_report.md & 复现指南与结果报告 \\
\bottomrule
\end{tabularx}
\caption{test13 相关主体结构分层（数据来源: 全库索引实测）}
\end{table}

% ================= 3 各部分关系 =================
\section{各部分关系}
\textcolor{relation}{依赖主链:} run-local.sh → main.py 子命令 → MultiGpuMultigrid →
CudaSchurOp ×N 线程（C++ libqcu.so 内核）→ nullvec\_cache 粗算子复用 → h5 产物归档
版本目录 v<ts>/。

\begin{keybox}
\textbf{数据流:} 每线程独立 params/argv/set\_ptrs（\_SET\_INDEX\_ 各自从 0 计数）；
求解热点在 worker 线程各自卡上并行（qcu.pyx with nogil 真并行）；多线程墙钟 =
max(各线程时间)。粗算子构建为 setup 阶段（缓存命中秒级），主线程 V100 完成。
\end{keybox}

\subsection{本次修复的三处读取 bug}
collect/mktable/plots/check 共用 \texttt{load\_dict\_h5} 还原结果，而
\texttt{save\_dict\_h5} 将 entries 列表写为数字 key 子组、读取时还原为 list；
原读取方按 dict 的 \texttt{.values()} 处理导致 AttributeError。修复为统一
\texttt{\_entries\_list()} 展开（main.py:557），并新增 \texttt{\_lat\_str()}
（main.py:622）兼容 dataset 还原键 \texttt{d\_lattice}；budget 子命令的
\texttt{--lattices} 解析补 levels=2（main.py:806）。

% ================= 4 项目思路 =================
\section{项目思路}
\subsection{设计意图}
多线程测试的核心动机：单线程 MG 已收敛（test12），需要验证"一线程一卡"在
P100×2 上的真并行加速；基准取多线程 BiStabCG 墙钟而非单线程，衡量的是
"并行 MG 相对并行 Krylov 求解器"的加速。数据持久化全部 h5py（独立 File
句柄，多线程安全），避免 json 并发写冲突。

\subsection{演进脉络}
test11（dev74 整合）→ test12（单文件优化版）→ test13（多线程版）；本次会话在
test13 上完成收尾：verify 重跑（新代码确认）→ bench 缩小 V100 组格子
（16x16x16x32 求解偏慢且缓存不全，超 30min 超时，改为 8x16x16x16）→
sweep/check/collect/mktable/plots/budget 全流程 → 报告。

\begin{ideabox}
\textbf{一句总结:} test13 验证了多线程 CUDA C++ MultiGrid 在小格子上的真并行
加速（2.1 倍），并以 h5py 全链路持久化保证多线程安全与跨环境可复现。
\end{ideabox}

% ================= 5 主题分析 =================
\section{主题分析}
\subsection{正确性验证（verify）}
全部 PASS：consistency（P100×2, 8x8x8x16, 2L, r5, rel\_diff=0.0）、
independent（4x4x4x8，两线程解不同且一致）、V100 单线程（8x16x16x16, 3L,
rel\_diff=0.0）、h5py 4 线程 IO 往返零误差。\textcolor{struct}{确证}——本次会话
实测输出（verify 命令 RC=0，PASS 行 4 条）。

\subsection{性能基准（bench）}
\begin{table}[htbp]\centering
\begin{tabular}{lrrrr}
\toprule
配置 & nT & ref(s) & mg(s) & speedup \\
\midrule
P100x2 8x8x8x16 2L & 2 & 0.571 & 0.272 & \textbf{2.101} \\
P100x2 8x8x8x16 3L & 2 & 0.575 & 0.451 & 1.275 \\
P100x2 8x16x16x16 2L & 2 & 0.613 & 1.024 & 0.598 \\
P100x2 8x16x16x16 3L & 2 & 0.598 & 1.495 & 0.400 \\
P100x2 16x16x16x16 2L & 2 & 0.695 & 0.779 & 0.892 \\
V100 16x16x16x16 3L & 1 & 0.679 & 0.734 & 0.926 \\
V100 8x16x16x16 3L & 1 & 0.600 & 0.528 & 1.137 \\
\bottomrule
\end{tabular}
\caption{bench 全 7 项（pairs=3；数据来源: test13\_bench.h5 实测）}
\end{table}

小格子 2L 加速显著（2.10）；3L 反而下降（粗层迭代开销占比上升）；
中格子 $<1$ 为已知历史特性（coarse solve 开销，AGENTS.md「已知硬件特性」）。

\subsection{参数扫描（sweep）与 gate}
16 配置全 PASS；gate=1.5 通过 11/16，最佳 2.149（L3 r10 ct1e5 cmi15）。
因子均值：2L=1.874 vs 3L=1.357；r10=1.732 vs r5=1.499；cmi15 优于 cmi10；
ct 影响弱。\textcolor{struct}{确证}——sweep 输出 16 行 OK + check 输出。

% ================= 6 关键代码片段 =================
\section{关键代码片段}
\subsection{粗算子构建统一 C++ matvec 路径}
\texttt{pyqcu/cuda/\_multi\_gpu.py:448-461}（未提交改动，本次 verify 确认正确）：
\begin{lstlisting}
        # 粗算子构建统一走 C++ matvec 路径（每线程一个 op；单线程 nthreads=1
        # 也用 1 个 CudaSchurOp，避免 Python matvec 构建大格子 50min+ 瓶颈
        ops_build = [CudaSchurOp(av, g, ce, coo, cei, coi, params=params_t)
                     for _ in range(max(1, min(self.nthreads, 4)))]
        coarse = build_schur_levels(
            op, S, self.num_levels, self.dof_list, self.mg_grid, self.lat_size,
            self.dof_list[1], self.dt, main_dev, nv_iters=self.nv_iters,
            use_cache=self.use_cache, cache_dir=self.cache_dir, verbose=False,
            matvec_ops=ops_build, nthreads=len(ops_build),
            av=av, params_template=params_t)
        for o in ops_build:
            o.release()
\end{lstlisting}

\subsection{entries 读取统一展开}
\texttt{logs/test13/main.py:557-574}（本次新增）：
\begin{lstlisting}
def _entries_list(d):
    """save_dict_h5 的 entries 组还原：{'entries': [{'e0': e}, ...]} →
    展开为 [e, ...]；兼容旧 dict 格式 {'e0': e, ...}。"""
    entries = d['entries']
    if isinstance(entries, dict):
        return [entries[k] for k in sorted(entries)]
    out = []
    for x in entries:
        if isinstance(x, dict):
            out.extend(v for v in x.values())
        else:
            out.append(x)
    return out
\end{lstlisting}

% ================= 7 参考源清单 =================
\section{参考源清单}
\begin{xltabular}{\textwidth}{llX}
\toprule
\textcolor{evidence}{参考源} & 类型 & 内容/作用 \\
\midrule
\endhead
\texttt{\nolinkurl{logs/test13/main.py:409-429}} & 仓库 & \_bench\_configs：bench 配置表（V100 组已缩格子） \\
\texttt{\nolinkurl{logs/test13/main.py:557-574}} & 仓库 & \_entries\_list：entries 统一展开（本次修复） \\
\texttt{\nolinkurl{logs/test13/main.py:622}} & 仓库 & \_lat\_str：兼容 d\_lattice 键（本次修复） \\
\texttt{\nolinkurl{logs/test13/main.py:806-808}} & 仓库 & budget --lattices 解析补 levels=2（本次修复） \\
\texttt{\nolinkurl{pyqcu/cuda/_multi_gpu.py:448-461}} & 仓库 & 粗算子构建统一 C++ matvec（未提交） \\
\texttt{\nolinkurl{logs/test13/test13_report.md}} & 仓库 & 本次收尾的结果报告（本报告数据源） \\
\texttt{\nolinkurl{logs/test13/AGENTS.md}} & 仓库 & 套件约定：参数/加速比语义/硬件特性 \\
\bottomrule
\end{xltabular}

% ================= 8 结论与局限 =================
\section{结论与局限}
\begin{conclbox}
\textbf{结论:} ① 结构——test13 为单文件套件 + 版本化产物目录，与 test12 同构；
② 关系——main.py 子命令 → MultiGpuMultigrid → CudaSchurOp×N 线程 → nullvec
缓存 → h5 归档，数据流清晰且多线程安全；③ 思路——真多线程并行验证达成：
小格子 2L 加速 2.10 倍，gate=1.5 通过 11/16，最佳 2.149；正确性全 PASS。
\end{conclbox}

\begin{warnbox}
\textcolor{warn}{局限:} ① 中/大格子（8x16x16x16、16x16x16x16）speedup $<1$
为已知特性，本次未做优化尝试（任务限定验证+参数收敛）；② 16x16x16x32
大格子 bench 未完成（求解慢 + 粗算子构建超 30min），已从配置表移除；
③ \texttt{\_multi\_gpu.py} 与 main.py 改动未提交（提示用户自行 git add/commit）；
④ 16x16x16x32 残留 2.8GB lv1 缓存（lv2 缺失不可用）待清理。
\end{warnbox}

\end{document}
